\section{Background and Motivating Example}
\label{sec:background}

\subsection{Dynamic Typing and Mutable State in Python}
Python associates a type with each object, while a name can be rebound to
objects of different types during execution~\cite{pythonExecutionModel,pythonDataModel}.
Attribute assignment similarly changes the value stored at an object
location. When several references designate the same object, an update
through one reference affects subsequent accesses through the others.
A call can therefore change the types encountered by its caller through
both its return value and updates to shared objects or global bindings.

Duck typing determines whether an operation is applicable from the behavior
supported by its operands~\cite{pythonGlossary}. The runtime objects and
their classes determine method lookup and operator dispatch. Different classes can
support the same operation with different result types and state effects.
For example, two encoder objects can both provide an \texttt{encode}
method while returning a string and bytes, respectively. Type annotations
declare interfaces for static checking~\cite{pep484}; an annotation allowing
both return types or field types still leaves the actual type at a
particular use dependent on the preceding execution.

We use \emph{type state} for the type information about relevant bindings
and object locations at a program point, together with the relations needed
to judge an operation. These relations include branch conditions and object
identity. A \emph{type-state transition} describes the change caused by a
statement or call. A function's \emph{type-state relation} connects its
input state and execution conditions to its return value and caller-visible
updates. Flow sensitivity tracks changes between program points; path
sensitivity retains the association between conditions and the states they
establish. If \texttt{prepare(obj, binary)} writes bytes when \texttt{binary}
is true and a string otherwise, a later string operation under
\texttt{not binary} uses the effect associated with that condition.

\subsection{Hoare Triples and State Relations}
\label{sec:hoare-background}
Hoare logic expresses a program's behavior using assertions about its
state~\cite{hoare1969axiomatic,softwareFoundationsHoare}. A triple
\begin{equation}
\{P\}\ C\ \{Q\}
\label{eq:hoare-triple}
\end{equation}
contains a precondition $P$, a command $C$, and a postcondition $Q$.
Under partial correctness, every execution of $C$ that starts in a state
satisfying $P$ and terminates normally ends in a state satisfying $Q$.
The judgment does not require termination. Assertions can describe types,
values, object identities, and relations between locations. For example,
with $o$ denoting the object referenced by \texttt{msg}, an encoding command
can relate a string in $\operatorname{Field}(o,\mathtt{body})$ before the
call to bytes in that same location afterward.

The composition principle connects commands through an intermediate
assertion: $\{P\}\ C_1\ \{R\}$ and $\{R\}\ C_2\ \{Q\}$ support
reasoning about $C_1;C_2$ from $P$ to $Q$. At a function call, this requires
instantiating the callee's assertions with the caller's arguments and object
identities. The resulting assertion supplies the state in which the next
operation is checked.

Our representation uses this assertion-based view to connect function
entries and exits, while retaining the intervening operations as an
ordered event graph. A case $(p,m,q)$ contains an entry applicability
assertion $p$, ordered behavior $m$, and a projected normal-exit assertion
$q$. The behavior records the checks and updates that establish the exit,
including prefixes ending in an operation failure. This makes intermediate
states and their supporting operations available to the detection and
refinement procedures in Section~\ref{sec:method}. Entry applicability
and internal operation requirements serve distinct roles: the former
selects a relation case, and the latter is checked against the state at
its execution point.

\subsection{Type-Error Categories and Analysis Entries}
An operation requirement follows from the semantics of the executed
operation, such as compatible operands for addition or a callable receiver
with valid argument binding. An annotation declares an interface against
which inferred behavior can be checked. Runtime type errors concern
operation failures; annotation-contract errors concern semantic
incompatibility between a declaration and behavior. Each category has its
own confirmation rule in Section~\ref{sec:conflict-analysis}.

The analysis takes declared entries. An open function entry supplies
symbolic arguments and relevant object and environment state; its input
requirements are inferred from function behavior. A repository entry
includes the initialization and calling sequence establishing that state.
Reports identify the entry under which a conflict is checked. The targeted
runtime failures involve incompatible operands, invalid argument binding,
non-callable objects, missing modeled attributes, and unsupported iteration
or indexing protocols. Exception propagation is part of the judgment:
a protocol failure handled within the entry is recorded as a handled
failure.

\subsection{A Call That Changes an Operation's Validity}
\label{sec:motivating-example}
Figure~\ref{fig:motivating-example} illustrates why type-error detection
requires tracking the state established by a call. The \texttt{Message}
class declares \texttt{body: str | bytes}. Both encoders accept a string. The \texttt{encode} method of the
text encoder returns that string; the UTF-8 encoder returns bytes. The helper
\texttt{encode\_message} writes the encoder's result back to the same
message object and implicitly returns \texttt{None}.

\begin{figure}[!htbp]
\centering
\includegraphics[width=\linewidth]{motivatingExample.pdf}
\caption{A call changes the type state used by a subsequent operation.
Point 1 establishes a string body; point 2 relates the selected encoder to
the field update; point 3 has a bytes body, making concatenation with a
string fail. The lower-right panel reports the illustrated example's
outcomes for Mypy, Pyrefly, Pyright, and our approach, PyTT.}
\Description{The caller make_line raises ValueError when the message body
is not a string. Otherwise it calls encode_message with Utf8Encoder, then
passes msg.body plus a newline string to print. The helper writes the
encoder result into msg.body. TextEncoder returns a string and Utf8Encoder
returns bytes. Numbered points mark the string state before the call,
the encoder-dependent update, and the bytes state after the call. The
result panel shows no diagnostic for Mypy, Pyrefly, and Pyright and a type
error detected by PyTT.}
\label{fig:motivating-example}
\end{figure}

\paragraph{From the guard to the update.}
In \texttt{make\_line}, a non-string body triggers \texttt{ValueError}.
Every execution reaching point~1 has therefore passed the guard with a
string body. The subsequent call supplies that message and a
\texttt{Utf8Encoder} instance to the helper. At point~2, the helper's two
encoder cases connect the input choice to different updates: the text
encoder establishes a string body and the UTF-8 encoder establishes bytes.
Both writes satisfy the field's union annotation. The relevant difference
is the state that each call establishes for subsequent operations.

\paragraph{From the update to the failure.}
Instantiating the UTF-8 case at the shown call binds the helper's message
to the caller's object. Consequently, point~3 contains a bytes body.
The string fact at point~1 concerns the old contents and does not constrain
the updated field. Evaluation of \texttt{msg.body + '\textbackslash n'}
then raises \texttt{TypeError} before \texttt{print} is invoked. With
\texttt{make\_line} as the analysis entry, that exception escapes the
function. Using \texttt{TextEncoder} at the same call would instead
establish a string body and permit the concatenation.

This reasoning distinguishes three representations. The helper's return
type, \texttt{None}, gives no description of the field update. A set of
possible field types, \texttt{str | bytes}, captures the possible contents.
A relation between the selected encoder, incoming body, and outgoing body
identifies the contents established in the actual calling context.
Preserving that relation lets the analyzer invalidate the old field fact
and check the addition using its new value.

The result panel in Figure~\ref{fig:motivating-example} identifies our
approach as \emph{PyTT}. For the illustrated example, it records no
diagnostic from Mypy, Pyrefly, or Pyright and a type error detected by PyTT;
the associated tool versions and configurations are xxx. The mechanism
illustrated by this result is the propagation of the encoder-dependent
write to the operation that consumes the field. Section~\ref{sec:evaluation}
evaluates detection across the benchmark under the stated configurations.

